
\documentclass{article}
\usepackage{colortbl}
\usepackage{makecell}
\usepackage{multirow}
\usepackage{supertabular}

\begin{document}

\newcounter{utterance}

\centering \large Interaction Transcript for game `bbh', experiment `snarks', episode 139 with reflexion{-}llama{-}3.1{-}8b{-}instruct.
\vspace{24pt}

{ \footnotesize  \setcounter{utterance}{1}
\setlength{\tabcolsep}{0pt}
\begin{supertabular}{c@{$\;$}|p{.15\linewidth}@{}p{.15\linewidth}p{.15\linewidth}p{.15\linewidth}p{.15\linewidth}p{.15\linewidth}}
    \# & \multicolumn{2}{c}{Player} && \multicolumn{2}{c}{Game Master} \\
    \hline

    \theutterance \stepcounter{utterance}  
    & & & \multicolumn{4}{p{0.6\linewidth}}{
        \cellcolor[rgb]{0.9,0.9,0.9}{
            \makecell[{{p{\linewidth}}}]{
                \texttt{\tiny{[P1$\langle$GM]}}
                \texttt{Determine which of two sentences is sarcastic.} \\
\\ 
\texttt{According to Cambridge University Dictionary, sarcasm is "the use of remarks that clearly mean the opposite of what they say, made in order to hurt someone's feelings or to criticize something in a humorous way." Sarcastic sentences often contain satirical or ironic utterances, hyperboles, ambivalent or witty remarks.} \\
\\ 
\texttt{Q: Which statement is sarcastic?} \\
\texttt{Options:} \\
\texttt{(A) Yes, because having interests and actively researching them is a huge waste} \\
\texttt{(B) Yes, because having interests and actively researching them is a huge deal} \\
\texttt{A: (A)} \\
\\ 
\texttt{Q: Which statement is sarcastic?} \\
\texttt{Options:} \\
\texttt{(A) No one is going to disagree with you on this. Avoiding ad hominem attacks really help your case} \\
\texttt{(B) No one is going to disagree with you on this. Ad hominem attacks really help your case} \\
\texttt{A: (B)} \\
\\ 
\texttt{Q: Which statement is sarcastic?} \\
\texttt{Options:} \\
\texttt{(A) Consistency in the league's punishments? What do you think this is supposed to be, politics?} \\
\texttt{(B) Consistency in the league's punishments? What do you think this is supposed to be, moral?} \\
\texttt{A: (A)} \\
\\ 
\texttt{Q: Which statement is sarcastic?} \\
\texttt{Options:} \\
\texttt{(A) Thank you stranger on the internet, I won't trust your three word post without any sources over someone offering solid advice claiming to actually be from the country in question} \\
\texttt{(B) Thank you stranger on the internet, I will trust your three word post without any sources over someone offering solid advice claiming to actually be from the country in question} \\
\texttt{A:} \\
            }
        }
    }
     \\ \\

    \theutterance \stepcounter{utterance}  
    & \multicolumn{4}{p{0.6\linewidth}}{
        \cellcolor[rgb]{0.561,0.737,0.561}{
            \makecell[{{p{\linewidth}}}]{
                \texttt{\tiny{[P1$\rangle$GM]}}
                \texttt{Based on the definition of sarcasm provided, I will analyze each option to determine which statement is sarcastic.} \\
\\ 
\texttt{1. Q: Which statement is sarcastic?} \\
\texttt{Options:} \\
\texttt{(A) Yes, because having interests and actively researching them is a huge waste} \\
\texttt{(B) Yes, because having interests and actively researching them is a huge deal} \\
\texttt{A: (A)} \\
\\ 
\texttt{This statement is sarcastic because it says the opposite of what it means. Having interests and actively researching them is generally considered a positive activity, not a waste. The use of "huge waste" is ironic, indicating that the speaker actually thinks it's a good thing.} \\
\\ 
\texttt{2. Q: Which statement is sarcastic?} \\
\texttt{Options:} \\
\texttt{(A) No one is going to disagree with you on this. Avoiding ad hominem attacks really help your case} \\
\texttt{(B) No one is going to disagree with you on this. Ad hominem attacks really help your case} \\
\texttt{A: (B)} \\
\\ 
\texttt{This statement is sarcastic because it says the opposite of what it means. Ad hominem attacks are generally considered to be unhelpful and counterproductive in an argument, not helpful. The use of "really help your case" is ironic, indicating that the speaker actually thinks ad hominem attacks are unhelpful.} \\
\\ 
\texttt{3. Q: Which statement is sarcastic?} \\
\texttt{Options:} \\
\texttt{(A) Consistency in the league's punishments? What do you think this is supposed to be, politics?} \\
\texttt{(B) Consistency in the league's punishments? What do you think this is supposed to be, moral?} \\
\texttt{A: (A)} \\
\\ 
\texttt{This statement is sarcastic because it says the opposite of what it means. The speaker is implying that the league's punishments are not consistent, and that's a bad thing. The use of "politics" is ironic, implying that the league's punishments are arbitrary and unfair, like in politics.} \\
\\ 
\texttt{4. Q: Which statement is sarcastic?} \\
\texttt{Options:} \\
\texttt{(A) Thank you stranger on the internet, I won't trust your three word post without any sources over someone offering solid advice claiming to actually be from the country in question} \\
\texttt{(B) Thank you stranger on the internet, I will trust your three word post without any sources over someone offering solid advice claiming to actually be from the country in question} \\
\texttt{A: (A)} \\
\\ 
\texttt{This statement is sarcastic because it says the opposite of what it means. The speaker is actually saying that they won't trust the stranger's advice, not that they will. The use of "I won't trust" is ironic, indicating that the speaker actually thinks the stranger's advice is untrustworthy.} \\
            }
        }
    }
    & & \\ \\

    \theutterance \stepcounter{utterance}  
    & & & \multicolumn{2}{p{0.3\linewidth}}{
        \cellcolor[rgb]{0.9,0.9,0.9}{
            \makecell[{{p{\linewidth}}}]{
                \texttt{\tiny{[GM$|$GM]}}
                \texttt{Target: (B)} \\
            }
        }
    }
    & & \\ \\

    \theutterance \stepcounter{utterance}  
    & & & \multicolumn{2}{p{0.3\linewidth}}{
        \cellcolor[rgb]{0.9,0.9,0.9}{
            \makecell[{{p{\linewidth}}}]{
                \texttt{\tiny{[GM$|$GM]}}
                \texttt{game\_result = LOSE} \\
            }
        }
    }
    & & \\ \\

    \theutterance \stepcounter{utterance}  
    & & \multicolumn{4}{p{0.6\linewidth}}{
        \cellcolor[rgb]{0.95,0.95,0.95}{
            \makecell[{{p{\linewidth}}}]{
                \texttt{\tiny{[P1$\leftrightarrow$P]}}
                \texttt{Post{-}episode Reflexion (not sent to GM):} \\
\texttt{New plan: } \\
\\ 
\texttt{Upon reflection, I realize that I was overly reliant on the definition of sarcasm provided and didn't consider the context and tone of each statement. I should have taken a more nuanced approach, considering the speaker's intent and the potential for irony or hyperbole. } \\
\\ 
\texttt{In the future, I will pay closer attention to the language and tone used in each statement, looking for signs of irony, hyperbole, or other forms of sarcasm. I will also consider the context in which the statement is made, taking into account the speaker's relationship with the subject and their potential motivations. } \\
\\ 
\texttt{Specifically, I will make sure to:} \\
\\ 
\texttt{* Be more aware of the potential for irony or hyperbole in each statement} \\
\texttt{* Consider the speaker's intent and potential motivations} \\
\texttt{* Look for signs of sarcasm in the language and tone used} \\
\texttt{* Take a more nuanced approach, rather than simply relying on the definition of sarcasm} \\
\\ 
\texttt{By taking these steps, I hope to improve my ability to identify sarcastic statements and achieve a better outcome in the future.} \\
            }
        }
    }
    & \\ \\

\end{supertabular}
}

\end{document}
